\section{Lifting the parameter and specializing to characteristic
\texorpdfstring{$p$}{p}}\label{lift:section}

We return to the curve and parameter of Theorem~\ref{intro:main}.
The remaining step is to lift the marked curve and its parameter to
characteristic zero, apply Proposition~\ref{con:charzero}, and
specialize the resulting perverse eigensheaf.  We retain the full
inertia homomorphisms in the lift, including when some $N_i$ vanish.

Let $R=W(k)$.  This is an excellent complete strictly henselian
discrete valuation ring, and $\ell$ is invertible in $R$.  There is
a smooth projective lift with disjoint marked sections
\[
 (\mathscr X,\mathfrak x_1,\ldots,\mathfrak x_s)
                         \longrightarrow\Spec R
\]
of $(X,x_1,\ldots,x_s)$.  Indeed, the obstruction group for the
marked curve is $H^2(X,T_X(-D))=0$; an ample line bundle also lifts,
and formal algebraization gives the projective family.  Use the
split models of $G$ and $B$ over $R$, and put
\[
 \mathscr D=\sum_i\mathfrak x_i,\qquad
 \mathscr U=\mathscr X\setminus|\mathscr D|,\qquad
 K_0=\overline{\operatorname{Frac}(R)}.
\]

\begin{lemma}\label{lift:parameter}
The parameter $\rho$ lifts to a continuous parameter
$\rho_{K_0}:\pi_1^{\et}(\mathscr U_{K_0})\to\Gd(L)$ with the same
compact image.  After compatible choices of boundary fibers and
prime-to-$p$ roots of unity, its full inertia homomorphism at
$\mathfrak x_i$ is
\begin{equation}\label{lift:inertia}
 \rho_{K_0}(\gamma)=\exp\bigl(t_{\ell,i}(\gamma)N_i\bigr).
\end{equation}
In particular every other prime factor of characteristic-zero
inertia acts trivially.  Its lattice reductions extend lisse over
$\mathscr U$ and specialize to those of $\rho$, with the tensor
compatibilities of Proposition~\ref{found:hecke-specialization}.
\end{lemma}

\begin{proof}
We use the finite-cover lifting argument of
\cite[Section~8.5, Lemma~8.4]{CT}, applied at all the disjoint
markings.  Let $H=\rho(\pi_1^{\et}(U))$, choose a cofinal system
$Q_r=H/H_r$ as in Lemma~\ref{con:parameter}, and let $Y_r^\circ$
be the corresponding connected finite \'etale torsor on $U$.
At every marking its inertia image has $\ell$-power order, since
the given inertia homomorphism factors through $\mathbb Z_\ell$.
Choose $d_r=\ell^{e_r}$ divisible by all these orders at stage $r$,
with $d_r\mid d_{r+1}$.  The torsor extends to the proper root stack
\[
 X_r=X\bigl(\sqrt[d_r]{x_1},\ldots,\sqrt[d_r]{x_s}\bigr).
\]
This is the root-stack description of tame covers
\cite[Proposition~A.10]{LieblichOlsson}.  The corresponding relative
root stack
\[
 \mathscr X_r=
 \mathscr X\bigl(\sqrt[d_r]{\mathfrak x_1},\ldots,
                         \sqrt[d_r]{\mathfrak x_s}\bigr)
\]
is smooth, proper, and tame with finite diagonal over $R$, since
$d_r$ is invertible in $R$.

Proper henselian invariance lifts the finite \'etale torsors to
$\mathscr X_r$ \cite[Theorem~4.5]{Rydh}; full faithfulness lifts
their group actions, transition maps after pullback to divisible
root stacks, and all compatibility relations.  See also
\cite[Theorem~A.11 and Corollaries~A.12--A.13]{LieblichOlsson}.
Restricting to $\mathscr U$ and then to its geometric generic fiber,
and choosing compatible base points, produces
\[
 \rho_{K_0}:\pi_1^{\et}(\mathscr U_{K_0})\longrightarrow H.
\]
Only the root indices must be prime to $p$; the finite group orders
$|Q_r|$ need not be.

The special torsor on $X_r$ is connected: its inverse image of $U$
is dense and is the connected torsor $Y_r^\circ$.  It is smooth and
geometrically reduced.  The proper flat tame-coarse-space argument
in Lemma~\ref{con:parameter} gives connected geometric generic
root-stack torsors.  A smooth connected curve stack is irreducible;
removing the inverse images of the finitely many marked gerbes
therefore leaves it connected.  Thus $\rho_{K_0}$ surjects onto
every $Q_r$, and its compact image is all of $H$.  In particular
the lifted parameter is Zariski dense.

To determine inertia, trivialize the normal line of each marked
section over the strictly henselian trait.  Its root gerbes are
then $B\mu_{d_r}$, with compatible atlases as $r$ varies.
Restricting the lifted torsor to this gerbe and pulling to its atlas
gives a constant finite torsor together with its $\mu_{d_r}$-action.
Full faithfulness of lifting retains exactly the special-fiber
action.  Compatible prime-to-$p$ roots of unity identify the generic
and special generators.  One may choose the boundary frames
compatibly throughout the tower: the possible frames form an
inverse system of nonempty finite fibers with surjective transition
maps.  At every stage, generic inertia consequently
factors through $\mu_{d_r}$ and agrees with the original inertia map.
Passing to the inverse limit proves \eqref{lift:inertia} for the
whole inertia group.  In particular its $\mathbb Z_a(1)$ factor
acts trivially for every prime $a\ne\ell$, including $a=p$.
This records the full peripheral compatibility, rather than only
the conjugacy class of one monodromy element; compare
\cite[Lemma~3.5]{LieblichOlsson}.

Finally, each finite reduction of an $H$-stable lattice in an
algebraic representation factors through some $Q_r$.  The lifted
torsor restricted to $\mathscr U$ is lisse and has the required
special fiber.  Tensor products and maps between the resulting
systems come from the same tower, with denominators cleared when
comparing lattices.  This proves the asserted lattice extension
condition.
\end{proof}

We may identify $K_0$ abstractly with $\mathbb C$.  Indeed $k$ is
countable and $W(k)$ is a set of sequences in $k$, while it contains
$\mathbb Z_p$.  Its cardinality, and hence that of $K_0$, is the
continuum.  An uncountable algebraically closed characteristic-zero
field has transcendence degree equal to its cardinality, so $K_0$
and $\mathbb C$ are isomorphic.  This identification is only of the
geometric base fields; the adic topology of the coefficient field
and continuity of the parameter remain unchanged.

\begin{proof}[Proof of Theorem~\ref{intro:main}]\label{lift:main-proof}
By Lemma~\ref{lift:parameter} and Proposition~\ref{con:charzero},
the ordinary-flag bundle stack over $K_0$ has a nonzero locally
constructible perverse eigenobject $M_{K_0}$ with eigenvalue
$\rho_{K_0}$.  Put
\[
 \mathscr A=\Bun_{G,B,\mathscr D}(\mathscr X/R),\qquad
                  M=\Psi M_{K_0}\quad\text{on }\mathscr A_k=\cA.
\]
The stack $\mathscr A$ is smooth over $R$: the unlevelled bundle
stack is smooth by vanishing of the degree-two obstruction groups,
and its flag fibers are products of the smooth variety $G/B$.
Proposition~\ref{found:nearby} implies that $M$ is locally
constructible and perverse.  Thus for a smooth finite-type chart
$f:S\to\cA$ of relative dimension $d$, the complex $f^*M[d]$ is
bounded constructible and perverse.  These are geometric nearby
cycles, without taking inertia invariants.

The Hecke geometry on $\mathscr U$ is the split smooth-leg geometry
of Section~\ref{found:section}.  The lattice extensions in
Lemma~\ref{lift:parameter} and
Proposition~\ref{found:hecke-specialization} give
\begin{equation}\label{lift:final-eigenmaps}
 \Hecke_{I,(V_i)}(M)\simeq
 M\boxtimes\Bigl(\boxtimes_{i\in I}(V_i)_\rho\Bigr)
 \qquad\text{on }\cA\times U^I.
\end{equation}
This is the relative Satake normalization
$E[d_\lambda](d_\lambda/2)$, with no moving-leg shift.  The
isomorphisms are natural in all $V_i$ and retain the unit,
convolution, permutation, and fusion compatibilities on every
collision diagonal, since specialization transports the coherent
system itself.

It remains to show that $M$ is nonzero.  Local constructibility gives
a $K_0$-valued point $(P,(\beta_i))$ of the geometric generic bundle
stack at which $M_{K_0}$ has nonzero stalk.  Choose a reference
bundle $P_0$ on $\mathscr X$, for example the trivial bundle.
A point of $U$ lifts to a section $y$ of $\mathscr U$ by
smoothness and henselianity.  On the affine curve
$\mathscr X_{K_0}\setminus\{y_{K_0}\}$, both $P$ and $(P_0)_{K_0}$
are trivial by the determinant argument used in
Section~\ref{con:section}.  Hence $P$ is a modification of
$(P_0)_{K_0}$ at $y_{K_0}$ in some finite Schubert bound.  After
a finite trait extension defining the data, properness of the
bounded unlevelled Hecke space extends the modification and thus
the bundle.  Each flag $\beta_i$ then extends by properness of the
associated $G/B$-bundle over the marked section.  We have extended
the original point to the relative ordinary-flag stack.

Lift this section to a finite-type smooth chart through its special
value, after a further henselian lift if necessary.  Its generic
stalk is still nonzero, so the closure of the generic nonzero-stalk
locus meets the special fiber on that chart.  Theorem~\ref{det:specialization}
forces $\Psi M_{K_0}\ne0$.  Its hypotheses hold by the Hecke
geometry above, the lattice extensions of
Lemma~\ref{lift:parameter}, and the special-fiber geometry of
Section~\ref{geom:section}.  Geometric nearby cycles are unchanged
by the finite trait extensions used in this argument.

Finally, Theorem~\ref{det:field}, applied to
\eqref{lift:final-eigenmaps}, gives $\SS(M)\subset\Lambda$ on every
smooth chart.  At ordinary flag level,
Definition~\ref{geom:cone} and its cotangent description identify
this with the cone of generically nilpotent Higgs fields
\[
 \phi\in H^0\bigl(X,\ad(P)\otimes\omega_X(D)\bigr),\qquad
 \Res_{x_i}\phi\in\Lie R_u(B_{\beta_i})\quad(1\leq i\leq s).
\]
The condition $p>n$ gives precisely the characteristic hypotheses
used in that geometry and detection.  The nonzero sheaf $M$ has
all the asserted properties.  Every parameter and sheaf used in
the construction is geometric; no Frobenius structure is required.
\end{proof}
